% TOP_PROBLEM: {"id": 1179, "title": "Schinzel's Hypothesis H", "category_id": 1, "category": {"id": 1, "name": "number_theory", "display_name": "Number Theory", "description": "Properties of integers, prime numbers, Diophantine equations.", "slug": "number-theory", "order_index": 1, "created_at": "2026-07-31T15:26:25.670Z"}, "status": "open"}
% FIRST_POSTED: null
% ENTRY_KIND: statement_only
% !TeX program = lualatex
% Definitions and sources retained from the original research notebook.
\documentclass[11pt]{article}
\usepackage[margin=25mm]{geometry}
\usepackage{fontspec}
\setmainfont{Latin Modern Roman}
\usepackage{amsmath,amssymb,mathtools}
\usepackage{unicode-math}
\setmathfont{Latin Modern Math}
\usepackage{enumitem,needspace}
\usepackage{xurl,hyperref}
\hypersetup{colorlinks=true,urlcolor=blue}
\setcounter{secnumdepth}{1}
\setlength{\parindent}{0pt}
\setlength{\parskip}{5pt}
\setlength{\emergencystretch}{3em}
\setlist[itemize]{leftmargin=3em,itemsep=5pt,topsep=4pt,parsep=0pt}
\allowdisplaybreaks
\newcommand{\N}{\mathbb{N}}
\newcommand{\Z}{\mathbb{Z}}
\newcommand{\Q}{\mathbb{Q}}
\newcommand{\R}{\mathbb{R}}
\newcommand{\C}{\mathbb{C}}
\newcommand{\F}{\mathbb{F}}
\newcommand{\A}{\mathbb{A}}
\newcommand{\PP}{\mathbb P}
\newcommand{\E}{\mathbb{E}}
\newcommand{\eps}{\varepsilon}
\newcommand{\dd}{\,\mathrm d}
\newcommand{\sref}[2]{\hyperlink{source:#1:#2}{[S#2]}}
\newcommand{\eref}[2]{\hyperlink{extra:#1:#2}{[E#2]}}
% Turn the next switch off for a shorter reading copy. The quotations preserve
% every exactTarget field, including qualifications omitted from a short summary.
\newif\ifshowcatalogue
\showcataloguetrue
\newcommand{\cataloguescope}[1]{\ifshowcatalogue
 \paragraph{Exact scope in the supplied catalogue.}
 \begin{quote}\small #1\end{quote}\fi}
\begin{document}
\setcounter{section}{0}
% ================================================================
% problemId: 1179
% Canonical problem ID: 1179
% exactTarget SHA-256: f079951f6ce8970223c01331fea5a722075ed301c3a6eacabd9032382ca38769
\clearpage
\section*{\texorpdfstring{Schinzel's Hypothesis H}{Schinzel's Hypothesis H}}\label{1179}
\subsection{Definitions and mathematical statement}
Let $f_1,\ldots,f_r\in\Z[x]$ be nonconstant irreducible polynomials with positive leading coefficients. The local admissibility condition is
\[
 \forall p\text{ prime}\ \exists n\in\Z:\quad p\nmid\prod_{i=1}^r f_i(n).
\]
The proposed conclusion is that $\{n\ge1:f_i(n)\text{ is prime for every }i\}$ is infinite. Irreducibility is over $\Q$ for primitive integer polynomials; local admissibility excludes fixed prime divisors. No asymptotic frequency is included in this qualitative assertion.
\par\smallskip\textit{Formulation and concept references:} \sref{1179}{1}.
\cataloguescope{Let f1,...,fr in Z[x] be irreducible with positive leading coefficients and no prime dividing f1(n)...fr(n) for every n. Prove that infinitely many positive integers n make every fi(n) prime.}
\subsection{Short English statement}
If irreducible integer polynomials have no fixed local obstruction, must they simultaneously take prime values infinitely often?
\subsection{Sources}
\begin{itemize}\raggedright
\item[\textnormal{[S1]}]\hypertarget{source:1179:1}{} Federico Scavia, Arithmetic Geometry and Algebraic Groups, Conjecture 3.1 (2026).
\url{https://www.lms.ac.uk/sites/default/files/inline-files/Scavia%20-%20Arithmetic%20Geometry%20and%20Algebraic%20Groups.pdf}.
{\small\textit{Catalogue formulation source.}}
\item[\textnormal{[S2]}]\hypertarget{source:1179:2}{} On The Schinzel–Wójcik Problem Under Hypothesis H.
\url{https://www.cambridge.org/core/journals/mathematical-proceedings-of-the-cambridge-philosophical-society/article/abs/on-the-schinzelwojcik-problem-under-hypothesis-h/C75A7C818639ADA780DE2A47C8BCFA19}.
{\small\textit{Catalogue research/status reference; not a proof endorsement.}}
\item[\textnormal{[S3]}]\hypertarget{source:1179:3}{} From Golomb to Schinzel.
\url{https://www.preprints.org/manuscript/202505.0651}.
{\small\textit{Catalogue research/status reference; not a proof endorsement.}}
\item[\textnormal{[E1]}]\hypertarget{extra:1179:1}{} Alexei N.\ Skorobogatov and Efthymios Sofos, \emph{Schinzel Hypothesis on average and rational points}, Inventiones Mathematicae 231 (2023), 673--739; arXiv:2005.02998; DOI 10.1007/s00222-022-01153-6. Introduction and Theorems 1.1--1.2.
\url{https://arxiv.org/abs/2005.02998}.
{\small\textit{Added research source: Fixed-family formulation and coefficient-averaged existence of a simultaneous prime value; the quantifier limitation is retained. Consulted 24 September 2026.}}
\item[\textnormal{[E2]}]\hypertarget{extra:1179:2}{} Tim Browning, Efthymios Sofos and Joni Ter\"av\"ainen, \emph{Bateman--Horn, polynomial Chowla and the Hasse principle with probability 1}, arXiv:2212.10373v2, 21 May 2026, Theorems 1.2 and 1.4 and the introduction.
\url{https://arxiv.org/abs/2212.10373v2}.
{\small\textit{Added research source: Averaged asymptotics and logarithmic exceptional-proportion savings in coefficient boxes; no result along every fixed affine orbit is assumed. Consulted 24 September 2026.}}
\item[\textnormal{[E3]}]\hypertarget{extra:1179:3}{} John Friedlander and Henryk Iwaniec, \emph{Asymptotic sieve for primes}, Annals of Mathematics 148 (1998), 1041--1065; arXiv:math/9811186, Section 1.
\url{https://arxiv.org/abs/math/9811186}.
{\small\textit{Added research source: Parity obstruction and the need for additional analytic information beyond ordinary congruence counts. Consulted 24 September 2026.}}
\item[\textnormal{[E4]}]\hypertarget{extra:1179:4}{} Huan Xiao, \emph{From Golomb to Schinzel}, Preprints.org manuscript 202505.0651, version 2, posted 9 June 2025, Section 3, especially the transition from (3.8) to (3.9) and the paragraph following (3.14).
\url{https://www.preprints.org/manuscript/202505.0651}.
{\small\textit{Added research source: Claimed solution, not endorsed. The attempt gives a coefficient-level counterexample to the displayed identity and an independent test of the limiting justification. Consulted 24 September 2026.}}
% END RESEARCH ADDITION REFERENCES-54
\end{itemize}

\end{document}
